\section{Industry Collaboration and Path Forward}

To realize this architecture requires multiple points of collaboration with the industry moving forward.

\subsection*{Align on Common Voltage Ranges, Connectors, and Current Levels}

Establishing standardized voltage levels, current ratings, and connector interfaces is essential to enabling modular and interoperable solutions across vendors and system integrators. For AI Factories adopting 800 VDC architectures, industry-wide alignment will minimize customization, accelerate time-to-deploy, and support scalability. A consistent framework allows for flexible deployment strategies while reducing complexity in design, installation, and maintenance.

\subsection*{Develop and Certify DC-Native Equipment}

The transition to facility-level DC distribution requires a new generation of DC-native equipment including rectifiers, distribution systems, protective devices, cabling, connectors, transfer trip schemes and metering. Coordinated development and certification efforts are needed to meet performance, safety, and reliability expectations for data center environments. Collaboration with equipment manufacturers and testing agencies is critical to expedite validation processes and build confidence in large-scale DC power adoption.

Key equipment categories requiring standardization:
\begin{itemize}
\item Medium-voltage to 800 VDC rectifiers (3.5--7.5 MVA range)
\item DC distribution boards with solid-state breakers
\item DC-optimized busducts and liquid-cooled cables
\item Touch-safe, mechanically-interlocked connectors for rack interfaces
\item DC-rated switchgear and protection devices
\item Integrated energy storage systems with fast-response controls
\end{itemize}

\subsection*{Align Safety Standards and Operational Practices}

Operating at 800 VDC voltages introduces new safety considerations that must be addressed through industry-aligned standards and procedures. Clear operational guidelines, including maintenance practices, arc flash mitigation, grounding strategies, and training, will ensure safe handling of 800 VDC infrastructure. Industry working groups and regulatory bodies must collaborate to update or expand applicable codes and best practices.

\subsection*{Priority Action Items}

\begin{enumerate}
\item \textbf{Form industry working groups} to define standardized voltage tolerances, connector pinouts, and current ratings for 800 VDC datacenter equipment.

\item \textbf{Accelerate certification programs} for DC-native components through UL, IEC, and other standards bodies.

\item \textbf{Develop reference implementations} for different scale points (1 MW, 10 MW, 100 MW) with documented bills of materials and integration guidelines.

\item \textbf{Establish testing facilities} for validating end-to-end system performance, fault response, and interoperability.

\item \textbf{Create training curricula} for electricians, maintenance staff, and operators on 800 VDC safety and operations.

\item \textbf{Coordinate with utilities} to define interconnection requirements and grid-support capabilities for AI Factory loads.
\end{enumerate}

\section*{Timeline for Adoption}

\begin{description}
\item[Near-term (0--2 years)] Deploy 800 VDC side-power racks and facility rectifiers using existing 1000 V-class components. Focus on pilot deployments and lessons learned.

\item[Mid-term (2--5 years)] Mature DC-native equipment ecosystem, complete certification processes, and establish industry standards. Begin transitioning from side-power to full facility-level 800 VDC distribution.

\item[Long-term (5+ years)] Evaluate migration to 1500 VDC distribution as safety standards solidify and component availability improves. Widespread SST deployment for direct MV-to-800 VDC conversion.
\end{description}

The window for industry alignment is narrow: AI Factory deployments are already underway, and early adopters will define the de facto standards. Proactive collaboration is essential to avoid fragmentation and ensure that the next generation of power infrastructure can keep pace with AI workload demands.
